\chapter*{KẾT LUẬN}
\addcontentsline{toc}{chapter}{KẾT LUẬN}

Đề tài đã xây dựng và đánh giá ba phương pháp phát hiện và định vị bất thường không giám sát trên đủ 15 loại của
MVTec AD theo cùng một giao thức. Autoencoder do nhóm tự thiết kế và huấn luyện; PaDiM và PatchCore dùng đặc trưng
học trước qua thư viện Anomalib.

Kết quả cho thấy PatchCore tốt nhất (image-AUROC 0,983, pixel-AUROC 0,980), sát số liệu của bài báo gốc. PaDiM đứng
thứ hai (0,890 và 0,960), nhanh nhất, nhưng điểm ảnh kém tin cậy ở các loại vật thể xoay. Autoencoder thấp nhất
(0,756 và 0,843): bản tái tạo mờ làm sai số xuất hiện ở viền và nền thay vì ở vùng lỗi, dù phương pháp vẫn tốt ở các
loại có lỗi lớn như toothbrush. Trong phạm vi thực nghiệm này, dùng lại đặc trưng ImageNet hiệu quả hơn nhiều so
với học tái tạo ảnh từ đầu trên vài trăm ảnh tốt.

Các hướng phát triển xuất phát trực tiếp từ hạn chế ở Chương 3:
\begin{itemize}
  \item Chạy Autoencoder với nhiều seed và thử hàm mất mát SSIM~\cite{bergmannImprovingUnsupervisedDefect2019} để
        biết phần nào của khoảng cách là do cấu hình.
  \item Chạy PaDiM và PatchCore với cùng backbone để tách ảnh hưởng của cơ chế khỏi ảnh hưởng của backbone.
  \item Chọn ngưỡng quyết định từ một tập ảnh tốt tách riêng và đo tỷ lệ báo động giả, làm cơ sở cho việc dùng mô
        hình trong một hệ thống kiểm tra thực tế.
\end{itemize}
